\documentclass[12pt]{article}
\usepackage[utf8]{inputenc}
\usepackage[T1]{fontenc}
\usepackage[margin=1in]{geometry}
\usepackage{setspace}
\usepackage{xcolor}
\usepackage[colorlinks=true, urlcolor=blue!60!black, linkcolor=blue!60!black]{hyperref}
\setlength{\parindent}{0em}
\setlength{\parskip}{0.6em}
\renewcommand{\baselinestretch}{1.0}
\begin{document}
\begin{center}
{\LARGE \textit{Obergefell v. Hodges}, 576 U.S. 644 (2015)\par}
\vspace{0.5em}
{\large Case Brief --- Two-Page Edition\par}
\end{center}

\section{Memory Jogger}
By 5--4, the Court held that the Due Process and Equal Protection Clauses require every state to license same-sex marriages and recognize those performed elsewhere, rejecting a rigid history-and-tradition ceiling on fundamental rights (\href{https://www.courtlistener.com/opinion/2812209/obergefell-v-hodges/}{576 U.S. 644}).

\section{Detailed Case Facts}
Same-sex couples challenged marriage bans in Michigan, Kentucky, Ohio, and Tennessee. James Obergefell married his dying partner John Arthur on a Maryland tarmac, but Ohio would not list him as surviving spouse; April DeBoer and Jayne Rowse could not both adopt their children; and Tennessee refused to recognize Sergeant Ijpe DeKoe's New York marriage.

\section{Procedural History}
Every district court ruled for the plaintiffs; the Sixth Circuit reversed, holding itself bound by \textit{Baker v. Nelson}, and the Supreme Court granted certiorari to resolve the circuit split.

\section{Judicial Votes}
Reversed, 5--4. Kennedy wrote for Ginsburg, Breyer, Sotomayor, and Kagan; Roberts, Scalia, Thomas, and Alito each dissented.

\section{Holding}
The right to marry is fundamental, and same-sex couples may not be denied it; states must license and recognize same-sex marriages on equal terms, and \textit{Baker v. Nelson} is overruled. Fundamental rights are identified by ``reasoned judgment,'' not \textit{Glucksberg}'s narrow formula.

\section{Analysis of Opinions}
Kennedy, invoking Harlan's \textit{Poe} dissent and the ``right to marry in its comprehensive sense'' recognized in \textit{Loving}, \textit{Zablocki}, and \textit{Turner}, identified four principles---autonomy, the two-person union, protection of children, and marriage as a keystone of social order---that apply equally to same-sex couples, and relied on the ``synergy'' of liberty and equality; ``fundamental rights may not be submitted to vote.'' Roberts asked ``who decides,'' accused the majority of ``effectively overrul[ing] \textit{Glucksberg},'' and saw no limit against plural marriage. Scalia decried a ``judicial Putsch''; Thomas called substantive due process a ``dangerous fiction'' protecting only freedom from restraint; Alito found no deeply rooted tradition and feared believers would be vilified.

\section{Examples: Future Applications}
\textbf{Same-Side:} refusing to list a birth mother's wife on a child's birth certificate withholds marital benefits (\textit{Pavan v. Smith}). \textbf{Opposite-Side:} age and consanguinity limits survive, since the holding reached only ``two consenting adults.'' \textbf{Fence-Sitter:} three co-parenting adults seeking plural marriage, as Roberts warned.

\section{Critique}
Originalists argue the 1868 meaning required nothing of the kind and that ``reasoned judgment'' substitutes judicial values for history; after \textit{Dobbs}, Thomas urged reconsideration. Progressives such as Kenji Yoshino and Laurence Tribe praised its fusion of liberty and equality but noted it left no scrutiny tier for sexual orientation.

\section{Key Quotations}
\begin{quote}
``The identification and protection of fundamental rights is an enduring part of the judicial duty to interpret the Constitution. That responsibility, however, `has not been reduced to any formula.' \textit{Poe v. Ullman}, 367 U.S. 497, 542 (1961) (Harlan, J., dissenting). Rather, it requires courts to exercise reasoned judgment in identifying interests of the person so fundamental that the State must accord them its respect. [\dots] History and tradition guide and discipline this inquiry but do not set its outer boundaries.''
\end{quote}
\noindent --- Justice Kennedy, Opinion of the Court, 576 U.S. at 664.

\end{document}
